\documentclass[10pt,a4paper]{amsart}
\usepackage[margin=19mm]{geometry}
\usepackage{amsmath,amssymb,mathtools}
\usepackage[scr=rsfs,cal=euler]{mathalfa}
\usepackage{xfrac}
\usepackage[unicode,hidelinks,bookmarks=false]{hyperref}
\newcommand{\ie}{i.e.\ }
\newcommand{\cf}{cf.\ }
\newcommand{\resp}{respectively\ }
\newcommand{\Subsection}{\subsection}
\providecommand{\nobreakdash}{\nobreak\hbox{-}}
\setlength{\emergencystretch}{2em}
\multlinegap0pt
\usepackage{amssymb}
\usepackage{mathtools}
\usepackage[shortlabels]{enumitem}
\usepackage{xcolor}
\usepackage{tikz}
\usepackage{tcolorbox}
\usepackage{float}
\newtheorem{theorem}{Theorem}
\newcommand{\Z}{\mathbb Z}
\newcommand{\mult}[2]{\binom{#1}{#2}}
\title{Optimal entanglement-assisted source coding\\ under a balanced-difference promise}
\author{Julius A. Zeiss}
\date{Source: arXiv:2609.15757v1 (CC BY 4.0)}
\begin{document}
\maketitle

\noindent\emph{Source and licence.} Optimal entanglement-assisted source coding\\ under a balanced-difference promise. Version arXiv:2609.15757v1 (CC BY 4.0), distributed by arXiv under the Creative Commons Attribution 4.0 International licence, http://creativecommons.org/licenses/by/4.0/, as recorded in the arXiv licence metadata for this exact version and shown on the abstract page https://arxiv.org/abs/2609.15757v1. Full licence notice: \texttt{licenses/arxiv-2609.15757-CC-BY-4.0.txt}.\par\smallskip

\noindent\textit{Editorial scope.} Counted unit: the article's theorem thm:permanent (source0.tex lines 1505--1511). The article's complete original proof of it is delivered with it (source0.tex lines 1587--1814, one \texttt{proof} environment of 228 lines, which the article places away from the statement). No mathematics is written, repaired or re-derived: the statement and the proof are the article's own lines, in the article's own notation, and no retained line is substituted or re-flowed.  No further source-local context is needed: every cross-reference of every retained line resolves inside the two delivered spans.  Nothing was omitted from any retained span, so the package carries no omission record.  No retained line contains a citation, so the wrapper carries no bibliography and no bibliography entry.  No retained line contains a figure environment, an image-inclusion command or drawing code, so no image is delivered and the package declares no asset dependency.  Editorial formatting only, in addition: an AMS \texttt{amsart} preamble that restates the article's own declarations and macros used by the retained lines, namely the environments \texttt{theorem} on separate counters, together with the standard packages those lines need.  The retained environments print in this document as Theorem 1 in order of appearance, whereas the article prints the same items with the numbering of its own full document.  The complete native source of the article as distributed in the arXiv e-print for this exact version is bundled unchanged as \texttt{sources/210447/source0.tex}.
\begin{theorem}[Uniform coefficient bound]\label{thm:permanent}
For every nonmonochromatic composition \(r\) of \(n=q\ell\),
\begin{equation}\label{eq:permanent}
 [z^r]H_{q,\ell}(z)\leq
 \frac{1}{n-1}\mult n{r_0,\ldots,r_{q-1}}.
\end{equation}
\end{theorem}

\begin{proof}[Proof of Theorem~\ref{thm:permanent}]
Recall that \(q\geq2\), \(\ell\geq1\), \(n=q\ell\), and
\((q-1)\ell\) is even under our standing assumptions.

\emph{1. Assignments and swaps.}
Fix a nonmonochromatic composition \(r\) of \(n=q\ell\).
Let
\begin{equation}
 \begin{aligned}
 X&=\{1,\ldots,\ell\}\times\Z_q,\qquad |X|=n,\\
 W_r&=\{w:X\to\Z_q:\ |w^{-1}(\{c\})|=r_c
                    \text{ for every }c\in\Z_q\}.
 \end{aligned}
\end{equation}
For any assignment \(w:X\to\Z_q\), define its target map by
\begin{equation}
 F_w(b,i)=(b,i+w(b,i)).
\end{equation}
All additions and subtractions within a block are in \(\Z_q\).
Let \(E_r=\{w\in W_r:F_w\text{ is bijective}\}\) and
\(I_r=W_r\setminus E_r\) be the valid and invalid assignments,
respectively. Choosing a permutation \(\pi_b\) of \(\Z_q\)
in each block is equivalent to choosing an assignment with bijective
target map, via \(w(b,i)=\pi_b(i)-i\). The corresponding monomial
in the permanent expansion records the total letter counts, so
extracting \([z^r]\) restricts this bijection to \(E_r\).
Together with the multinomial count of all assignments, this gives
\begin{equation}\label{eq:count}
 |W_r|=\mult n{r_0,\ldots,r_{q-1}},\qquad
 |E_r|=[z^r]H_{q,\ell}(z).
\end{equation}
In particular, \(|W_r|>0\), and it suffices to prove
\(|E_r|/|W_r|\leq1/(n-1)\).

Write \(\binom X2\) for the set of unordered pairs of distinct
positions. For \(p=\{u,v\}\in\binom X2\), let \(\tau_p\)
exchange \(u,v\) and fix every other position, and put
\(w^p=w\circ\tau_p\). Thus \(w^p\in W_r\) and
\((w^p)^p=w\). We consider only pairs for which \(w(u)\ne w(v)\).
A \emph{repair} of \(w\in I_r\) is such a pair \(p\) with
\(w^p\in E_r\). Pairs may involve positions in different blocks.

\emph{2. The defects caused by a swap.}
For \(c\in\Z_q\), define the fixed shift
\begin{equation}
 T_c(b,i)=(b,i+c).
\end{equation}
This is a permutation of \(X\), and
\(F_w(u)=T_{w(u)}(u)\). For each fixed \(u\), the targets
\(T_c(u)\) are distinct as \(c\) varies.

Let \(w\in E_r\), \(p=\{u,v\}\), and
\(A=w(u)\ne w(v)=B\). Denote the old and new targets by
\begin{equation}
 \alpha=T_A(u),\quad \beta=T_B(v),\qquad
 \gamma=T_B(u),\quad \delta=T_A(v).
\end{equation}
Bijectivity of \(F_w\) gives \(\alpha\ne\beta\).
Moreover,
\(\{\alpha,\beta\}\cap\{\gamma,\delta\}=\varnothing\):
for example, \(\gamma\ne\alpha\) since \(A\ne B\), and
\(\gamma\ne\beta\) since \(T_B\) is injective and \(u\ne v\).
The two comparisons involving \(\delta\) follow in the same way.
Only the targets of \(u,v\) change. Hence \(\alpha,\beta\)
have no preimages under \(F_{w^p}\). Each distinct new target
\(\gamma\) or \(\delta\) already had a unique preimage under
the bijection \(F_w\). That preimage lies outside \(\{u,v\}\),
because the old targets of \(u,v\) were \(\alpha,\beta\), and
therefore remains unchanged by the swap.
If \(\gamma\ne\delta\), the swap adds one preimage to each,
so each has exactly two preimages. If \(\gamma=\delta\), the swap
adds both \(u\) and \(v\) as preimages of their common target,
which therefore has exactly three.
Every remaining target has exactly one preimage. In particular,
\(w^p\in I_r\).

It follows also that every invalid assignment admitting a repair
has exactly one of these two defect patterns: apply the preceding
argument to the valid assignment obtained by that repair and reverse
the swap. When two targets each have two preimages, any repair must
select one position from each of these two preimage sets. When one
target has three preimages, it must select two of those three
positions. Otherwise, an unchanged pair of positions would still
have the same target after the proposed repair.

\emph{3. The double-counting identity.}
Choose either all swaps of unequal letters or only swaps exchanging
a specified letter \(c\) with a different letter. For \(w\in W_r\),
let \(\mathcal P(w)\subseteq\binom X2\) be the pairs permitted
by the chosen rule. Both rules are preserved by reversing a swap:
\begin{equation}
 p\in\mathcal P(w)\quad\Longleftrightarrow\quad
 p\in\mathcal P(w^p).
\end{equation}
By Step~2, every permitted swap from \(E_r\) ends in \(I_r\).
The map \((w,p)\mapsto(w^p,p)\), whose inverse is the same map,
therefore bijects permitted swaps from valid assignments with
permitted repairs of invalid assignments. Consequently,
\begin{equation}\label{eq:switching-count}
 \sum_{w\in E_r}|\mathcal P(w)|
 =\sum_{v\in I_r}
   |\{p\in\mathcal P(v):v^p\in E_r\}|.
\end{equation}
We count swaps together with their unordered position pairs, so each
swap occurs once on each side of this identity. If every valid
assignment has \(d\) permitted swaps and every invalid assignment
has at most \(R>0\) permitted repairs, then
\begin{equation}\label{eq:switching-density}
 d|E_r|\leq R|I_r|=R(|W_r|-|E_r|),\qquad
 \frac{|E_r|}{|W_r|}\leq\frac{R}{d+R}.
\end{equation}

\emph{4. The composition cases.}
We distinguish the following three exhaustive cases.

\emph{Case 1: a letter occurs between \(2\) and \(n-2\) times.}
Suppose that \(r_c=k\) with \(2\leq k\leq n-2\).
Permit only swaps exchanging \(c\) with a different letter.
There are exactly \(d=k(n-k)\) permitted pairs for every assignment:
choose one of the \(k\) positions carrying \(c\) and one of the
\(n-k\) other positions.

Fix \(v\in I_r\). If it has no permitted repair, the required
bound is immediate. Otherwise, Step~2 applies. Each preimage set of
\(F_v\) contains at most one position carrying \(c\), because all
such positions are mapped by the injective map \(T_c\).
If two targets each have two preimages, let
\(\varepsilon,\delta\in\{0,1\}\) count the positions carrying
\(c\) in their respective preimage sets. A permitted repair must
choose one position from each set, exactly one of which carries
\(c\). The number of candidate pairs is therefore
\begin{equation}
 \varepsilon(2-\delta)+\delta(2-\varepsilon)\leq2.
\end{equation}
If one target has three preimages, a permitted repair must choose
two of them. At most one carries \(c\), leaving at most two
candidate pairs. Thus every invalid assignment has at most two
permitted repairs; no claim that every candidate pair succeeds is
needed. Taking \(R=2\) in~\eqref{eq:switching-density} gives
\begin{equation}\label{eq:two}
 \begin{aligned}
 k(n-k)|E_r|&\leq2\bigl(|W_r|-|E_r|\bigr),\\
 \frac{|E_r|}{|W_r|}&\leq\frac2{k(n-k)+2}.
 \end{aligned}
\end{equation}
Since
\begin{equation}
 k(n-k)-2(n-2)=(k-2)(n-k-2)\geq0,
\end{equation}
we obtain
\begin{equation}
 \frac{|E_r|}{|W_r|}
 \leq\frac2{2(n-2)+2}=\frac1{n-1}.
\end{equation}

\emph{Case 2: a letter occurs \(n-1\) times.}
Let this letter be \(c\), and take any \(w\in W_r\).
Subtracting \(c\) from every letter gives
\(F_{w-c}=T_{-c}\circ F_w\), so it preserves validity.
Since the composition is nonmonochromatic, \(w-c\) is zero at
all but one position \(u\), where it is nonzero. Thus
\(F_{w-c}\) fixes every position except \(u\). Its target
\(v=F_{w-c}(u)\ne u\) also satisfies \(F_{w-c}(v)=v\),
so \(F_{w-c}\) is not injective and \(w\) is invalid.
Hence \(E_r=\varnothing\), and the coefficient bound holds.

\emph{Case 3: neither of the preceding cases applies.}
Nonmonochromaticity excludes an entry \(r_c=n\), so every positive
entry of \(r\) must be one. Since there are only \(q\) letters,
this gives \(n\leq q\); combined with \(n=q\ell\geq q\), it
forces \(\ell=1\), \(n=q\), and
\(r=(1,\ldots,1)\). Since \((q-1)\ell\) is even, \(n=q\)
is odd. There is now one block, so we identify its positions with
\(\Z_n=\Z_q\).

Permit all swaps. Every assignment in \(W_r\) has distinct letters,
so it admits \(d=\binom n2\) permitted pairs. We prove that each
invalid assignment has at most three repairs. An assignment without
a repair already satisfies this bound. Otherwise, Step~2 applies.
If one target has three preimages, a repair must exchange letters
at two of them, giving at most \(\binom32=3\) choices.

For the other defect pattern, write \(U,V\in\Z_q\) for the
two targets with two preimages, whose respective preimage sets are
\(\{p_1,p_2\}\) and \(\{s_1,s_2\}\). Let \(H_1\ne H_2\)
be the two missing targets. A repair must exchange letters at some
\(p_i\) and \(s_j\), giving four candidate pairs. The two
preimage sets are disjoint because \(U\ne V\). Since the letters at these positions
are \(U-p_i\) and \(V-s_j\), their new targets after a swap are
\begin{equation}
 V+(p_i-s_j),\qquad U-(p_i-s_j).
\end{equation}
A repair must send both changed positions to the two missing targets;
in particular, a necessary condition is
\begin{equation}
 \delta_{ij}:=p_i-s_j\in
 \{H_1-V,H_2-V\}=\{h,k\},\qquad h\ne k.
\end{equation}
If all four pairs repaired the assignment, all four \(\delta_{ij}\)
would belong to this two-element set. Each row and each column of
\((\delta_{ij})\) has distinct entries, because the positions within
either pair are distinct. Up to exchanging \(h,k\), the matrix
would consequently be
\begin{equation}
 \begin{pmatrix}h&k\\k&h\end{pmatrix}.
\end{equation}
But
\(\delta_{11}+\delta_{22}=\delta_{12}+\delta_{21}\)
would then give \(2h=2k\) in \(\Z_q\).
Since \(q\) is odd, multiplication by two is injective, contradicting
\(h\ne k\). At most three of the four pairs can therefore repair the
assignment.

Applying~\eqref{eq:switching-density} with
\(d=\binom n2=n(n-1)/2\) and \(R=3\) yields
\begin{equation}\label{eq:three}
 \begin{aligned}
 d|E_r|&\leq3\bigl(|W_r|-|E_r|\bigr),\\
 \frac{|E_r|}{|W_r|}
 &\leq\frac3{d+3}
 =\frac3{n(n-1)/2+3}\leq\frac1{n-1}.
 \end{aligned}
\end{equation}
The last inequality is equivalent to
\((n-3)(n-4)\geq0\), which holds for every odd integer \(n\geq3\).
All compositions have now been covered, proving
\eqref{eq:permanent} by \eqref{eq:count}.
\end{proof}

\end{document}
